\documentclass[UTF8,zihao=-4]{ctexart}

\usepackage[a4paper,top=2.4cm,bottom=2.4cm,left=2.25cm,right=2.25cm]{geometry}
\usepackage{amsmath,amssymb}
\usepackage{booktabs,longtable,array,tabularx}
\usepackage{enumitem}
\usepackage{xcolor}
\usepackage{hyperref}
\usepackage{fancyhdr}

\definecolor{inkblue}{RGB}{34,76,112}
\definecolor{softgray}{RGB}{246,248,250}
\hypersetup{
  colorlinks=true,
  linkcolor=inkblue,
  urlcolor=inkblue,
  citecolor=inkblue,
  pdfauthor={数学建模学习小组},
  pdftitle={数学建模阶段学习笔记与论文写作模板}
}

\pagestyle{fancy}
\fancyhf{}
\fancyhead[L]{数学建模阶段学习笔记与论文写作模板}
\fancyhead[R]{\thepage}
\renewcommand{\headrulewidth}{0.4pt}
\setlength{\headheight}{15pt}
\setlength{\parindent}{2em}
\setlength{\parskip}{0.25em}
\setlist{nosep,leftmargin=2.2em}
\renewcommand{\arraystretch}{1.28}

\newenvironment{writeguide}
  {\begin{quote}\small\color{black}}
  {\end{quote}}
\newcommand{\guideitem}[2]{%
  \par\noindent\textbf{\color{inkblue}#1：}\ignorespaces #2\par
}
\newcommand{\keywords}[1]{\par\noindent\textbf{关键词：}#1}

\title{\bfseries 数学建模阶段学习笔记与论文写作模板}
\author{队长、组员1、组员2}
\date{2026年8月}

\begin{document}

\maketitle

\begin{abstract}
本报告记录本队现阶段的算法学习、论文精读和写作模板。第一部分依照思维导图，
按三人分工整理十三个学习单元，列出基本原理、求解步骤、代码入口和适用边界；
第二部分选取2024年A题一等奖论文A053《基于目标优化的板凳龙形态调节》，分析其
摘要、假设、建模、求解、结果解释与附录的写法；最后给出本队以后写论文时采用的
章节模板。文中的“已学习”指已能复述原理、说明步骤并列出代码逻辑，不等同于已经
完成全部算法的限时独立复现。

\keywords{数学建模；学习笔记；板凳龙；论文结构；写作模板}
\end{abstract}

\tableofcontents
\newpage

\section{阶段学习笔记}

\subsection{选题侧重与成员分工}

从近六年A、B、C题的任务差异看，A题常要把物理过程写成方程并处理数值误差，
B题更看重决策变量、约束和可行方案，C题则从数据口径、预测或评价进入决策。
据此，队长负责机理与连续模型，组员1负责规划和优化，组员2负责数据与统计学习。
三人共用同一套输入输出表：队长给出状态量和物理边界，组员1接收状态量建立目标
与约束，组员2整理参数、特征和误差。这样分工后，一项结果由谁计算、交给谁使用，
在代码和论文中都能找到对应位置。

本阶段按三项标准记录学习成果：能用自己的话说清模型为何适用；能把公式写成程序
步骤；能指出它在什么条件下不该使用。思维导图中的方法很多，本次先选十三个单元，
以后再用真题补足尚未动手的分支。

\subsection{十三个算法与模型学习单元}

\setlength{\tabcolsep}{2.5pt}
\small
\begin{longtable}{p{1.25cm}p{2.55cm}p{6.25cm}p{5.45cm}}
\caption{成员学习内容、代码逻辑与使用边界}\label{tab:learning-units}\\
\toprule
成员 & 算法或模型 & 基本原理、步骤与代码逻辑 & 适用问题、优缺点和限制\\
\midrule
\endfirsthead
\toprule
成员 & 算法或模型 & 基本原理、步骤与代码逻辑 & 适用问题、优缺点和限制\\
\midrule
\endhead
\bottomrule
\endfoot

队长 & ODE与RK4
& 从受力、守恒或状态转移写出 $\dot{y}=f(t,y)$，给定初值后按四个斜率加权推进。
代码入口为 \texttt{rhs(t,y,p)} 和 \texttt{rk4\_step(...)}；用Euler法和
\texttt{solve\_ivp}/\texttt{ode45}作对照，并保存步长误差。
& 适合状态随时间连续变化的力学、传热、种群和传播过程。RK4精度与实现难度较均衡，
但刚性方程会受稳定域限制；此时应改用刚性求解器，不靠缩小一步掩盖问题。\\

队长 & PDE与有限差分
& 写明空间区域、初始条件和边界条件，将时间、空间导数换成差分格式；代码依次生成
网格、装配差分矩阵、推进时间层，并比较三档网格的同一观测量。
& 适合规则区域上的热、波和扩散问题。公式与网格位置对应直观，代价是复杂边界难处理，
显式格式还须核对稳定条件。\\

队长 & 有限元模型
& 将PDE写成弱形式，把区域划分为单元；逐单元计算局部矩阵，装配全局稀疏矩阵，施加
边界后求节点量。代码模块依次为 \texttt{mesh}、\texttt{element}、\texttt{assemble}、
\texttt{bc} 和 \texttt{solve}。
& 适合不规则几何、变系数介质和混合边界。几何适应性强，但网格质量、形函数和边界
装配都容易引入错误；规则区域的入门题先用有限差分作基线。\\

队长 & 几何递推与事件判定
& 由已知龙头位置建立相邻把手的定长方程，筛去不满足构件次序和运动方向的根，再逐节
递推。碰撞用矩形顶点和分离轴判定，时间推进遇到第一次触发即停。代码保留
\texttt{next\_handle}、\texttt{collision} 和状态码。
& 适合链式刚体、轨迹跟随和首达事件。优点是公式、程序循环和物理对象能对读；缺点是
多根、步长和漏检风险要另列检查，不能只写“取合理解”。\\

组员1 & LP与MILP
& 写出决策变量、线性目标、等式和不等式；启停、选择或分配用0--1变量。代码先生成
小规模枚举答案，再调用PuLP/Pyomo或MATLAB求解器，读取可行性、目标值和最优间隙。
& 适合资源分配、排程、选址和生产决策。线性模型便于解释并能给最优性信息；整数规模
大时计算量上升，Big--M须来自实际上界，不能任意指定过大的常数。\\

组员1 & 多目标与NSGA--II
& 分别保留各目标，不先强行合成一个分数；染色体经过可行性修复后，按非支配排序和拥挤
距离选择，输出Pareto解集。代码记录种群、代数、随机种子和每个候选的原始指标。
& 适合成本、收益、风险等目标互相冲突的问题。它能展示取舍，但不替参赛者选最终方案；
变量少且约束凸时，$\varepsilon$约束法可直接呈现取舍边界。\\

组员1 & 粒子群优化
& 每个粒子保存位置、速度、个体最好解和群体最好解；迭代中更新速度、处理越界并重新
评价。代码把目标函数、约束罚项、种子和停止阈值分开，至少做多次独立运行。
& 适合连续、非光滑或难求导的参数搜索，编码简短；它没有连续全局最优证书，结果会受
边界、罚项和随机种子影响。若变量只有一维且单调，可改用遍历或二分。\\

组员1 & 图论与网络优化
& 把对象写成节点、关系写成带权边：非负单源最短路用Dijkstra，容量与费用同时存在时
用最小费用流。代码先建图，再运行算法，最后把路径或流量映射回现实对象并核对容量。
& 适合运输、路径、匹配和网络资源配置。结构清楚、结果易复算；动态图、负权边和多目标
路线要另行处理，最小生成树也不能代替两点间最短路。\\

组员2 & 缺失、异常与重复处理
& 先建数据字典和主键，区分未检出、未知和未采样；缺失值再按成因选择删除、插值或KNN，
异常值用箱线图、$3\sigma$和业务范围复核。代码以训练集拟合清洗参数并输出处理记录。
& 适合一切附件数据题。它决定后续样本是否可信，但不能把业务极值自动删除，也不能在
划分训练集前用全体数据估计均值或尺度。\\

组员2 & PCA与K--Means
& 标准化后由协方差矩阵特征分解得到主成分，按解释率选维数；再对样本运行K--Means，
比较不同簇数的轮廓系数。代码输出载荷、解释率、簇中心和原样本编号。
& PCA适合线性降维，K--Means适合近似球形簇。两者速度快、图形直观，但主成分未必有
直接业务含义，密度不均或非凸簇应比较DBSCAN或谱聚类。\\

组员2 & 熵权与TOPSIS
& 指标先正向化和无量纲化，熵权反映样本差异，TOPSIS再计算对象到正、负理想解的距离。
代码逐步导出标准矩阵、权重、距离和排序，并对权重作扰动。
& 适合多指标评价与方案排序。计算过程清楚，但“差异大”不等于“业务上更重要”；指标
方向、相关重复和权重敏感性都要说明，不能只给最终名次。\\

组员2 & 回归与分类基线
& 连续响应以线性、Ridge和Lasso为基线，类别响应以Logistic为基线；在同一训练/验证划分下再与
随机森林或XGBoost比较。代码用pipeline串起预处理、拟合、预测和指标，保存系数或特征贡献。
& 适合关系估计、风险评分和类别判断。线性模型可解释，树模型能处理非线性；后者容易
过拟合且外推较弱。相关或高准确率都不构成因果证据。\\

组员2 & ARIMA与滚动验证
& 先画序列并检查趋势、季节和差分后的平稳性，再选ARIMA/SARIMA阶数；预测时只用当前
时点以前的数据，按滚动窗口计算MAE、RMSE和MAPE。代码保存每个预测起点和误差。
& 适合有时间顺序的需求、价格和状态预测。统计模型对中小样本友好，但结构突变会使历史
规律失效；时间序列不能随机打乱做普通K折。\\
\end{longtable}
\normalsize

\subsection{本阶段做得较好的地方}

一是分工落到了模型接口，而不是只分配算法名称。比如队长输出位置和速度，组员1把它们
写进速度约束，组员2负责参数表和误差口径，三人的结果可以在同一张变量表中衔接。
二是每个学习单元都留有一个简单方法作比较：RK4对照Euler，MILP对照小规模枚举，
粒子群对照遍历，树模型对照线性或Logistic。三是论文学习没有停在目录层面；A053中
“定长递推--物理选根--第一次碰撞--回退求精--分区状态”的顺序已经整理成可照着检查
代码的笔记。

\subsection{目前不足}

现有材料能说明三人看过哪些内容，却还不能证明三人在比赛时限内都能独立写完代码。
队长对复杂边界的有限元网格调试不熟；组员1尚缺大规模MILP不可行诊断和随机算法多次
运行记录；组员2对时间序列泄漏、类别不平衡和权重扰动的动手样例偏少。全队也没有用
一套结果文件完整检查“代码数值--正文表格--摘要结论”是否一致。这些问题下一阶段用
真题产物检查，不用“继续加强学习”一笔带过。

\subsection{下一阶段四周安排}

\small
\begin{longtable}{p{1.3cm}p{1.3cm}p{3.9cm}p{4.7cm}p{4.15cm}}
\caption{四周训练任务与核对办法}\label{tab:four-week-plan}\\
\toprule
时间 & 负责人 & 任务 & 产物 & 核对办法\\
\midrule
\endfirsthead
\toprule
时间 & 负责人 & 任务 & 产物 & 核对办法\\
\midrule
\endhead
\bottomrule
\endfoot
第1周 & 队长 & 用一道A题完成ODE/RK4基线，并完成一个二维扩散FDM算例。
& 方程说明、两套脚本、三档步长或网格结果。 & 量纲正确；误差随加密下降；另一名成员可按入口复算一张表。\\
第1周 & 组员1 & 用小规模排程题对照枚举与MILP。 & 变量表、约束表、枚举结果和求解日志。 & 小样例目标值一致；每条约束能指出题意来源。\\
第1周 & 组员2 & 完成一份附件数据的字段、缺失、异常和泄漏检查。 & 数据字典、清洗记录、训练/验证划分脚本。 & 每次删除或填补都有对象和理由；预处理不读取验证集统计量。\\
第2周 & 组员1 & 在同一连续优化题上比较遍历、PSO和NSGA--II的对应任务。 & 参数配置、20个种子结果、收敛与可行率表。 & 计算预算一致；报告中位数和离散程度，不只报最好一次。\\
第2周 & 组员2 & 完成PCA--聚类和熵权--TOPSIS两个小例。 & 解释率、载荷、簇中心、权重扰动与排序表。 & 能说明主成分、簇和排名分别回答什么问题。\\
第3周 & 全队 & 选择一题做36小时压缩训练。 & 可编译论文、代码入口、结果目录和复盘单。 & 正文中每个关键数值可追到结果文件；无未定义引用。\\
第4周 & 全队 & 安排72小时完整训练并互换审稿。 & 论文、源代码、附件、个人问题清单。 & 12小时内定题，24小时给基线，48小时冻结主结果，交稿前完成复算。\\
\end{longtable}
\normalsize

\section{一等奖论文A053学习分析}

\subsection{选文与阅读范围}

本次选择2024年A题编号A053的《基于目标优化的板凳龙形态调节》\cite{A053}。
论文共62页，前42页为正文和参考材料，后20页为MATLAB代码。选择它不是因为模型名称多，
而是因为五个分问共用同一条几何递推主线：前一问求得的位置、碰撞对象和可行边界，
在后一问中继续充当输入。正文中的四类路径状态还能与附录状态码1--4对读，便于学习
论文怎样把公式、算法和代码连在一起。

\subsection{论文结构及各部分写法}

\small
\begin{longtable}{p{2.25cm}p{4.0cm}p{5.1cm}p{4.2cm}}
\caption{A053章节功能与本队理解}\label{tab:a053-structure}\\
\toprule
部分 & 论文在写什么 & 具体写法 & 本队吸收的做法\\
\midrule
\endfirsthead
\toprule
部分 & 论文在写什么 & 具体写法 & 本队吸收的做法\\
\midrule
\endhead
\bottomrule
\endfoot
摘要（p.1）
& 依次回答运动状态、碰撞、最小螺距、调头曲线和速度上限。
& 每个分问都给对象、模型动作和交付值，如第一次碰撞时刻、螺距和龙头速度；算法名放在它所求的量之后。
& 摘要按分问写，每问至少留下一个有判别力的结果；减少重复的顺序词，不把摘要写成章节目录。\\

问题分析（p.3--4）
& 把板凳龙的整体运动拆成相邻把手位置递推，再说明各问增加的约束。
& 问题四被拆成“求较短调头曲线”和“求调头过程中的运动状态”两个子任务，前者输出路径，后者沿用位置递推。
& 先画清输入输出关系，再选方法；复用前问时点明沿用量和新增量。\\

模型假设（p.2）
& 限定人为扰动、板凳刚体和开始调头的时机。
& 假设会进入方程或事件判定，没有写与计算无关的常识句。
& 每条假设说明“它使哪一项可以怎样计算”，并在评价中说明放宽后要改哪个模块。\\

符号说明（p.3）
& 给出螺距、极角、位置、速度、圆弧半径等贯穿五问的量。
& 统一公共符号，分问新增量在首次使用处定义，减少同一字母承担两个含义。
& 符号表只放多次使用的量；仅在一处使用的中间量随公式解释。\\

模型建立（p.6--21、p.23--36）
& 写螺线弧长、定长递推、矩形碰撞、相切圆弧和多段路径状态。
& 公式按“已知把手--定长方程--候选根--运动方向筛选--下一把手”推进；碰撞从几何对象落到矩形投影区间。
& 让公式顺序接近程序顺序。遇到多根时写清筛根条件，遇到分段路径时先列状态和切换条件。\\

模型求解（p.8--16、p.22--31、p.34--41）
& 用二分、时间离散、变步长遍历、状态递推和粒子群得到五问结果。
& 螺距先在 $[45,46]\,\mathrm{cm}$ 内定位，再以 $0.01\,\mathrm{cm}$ 步长细查；碰撞检测在第一次触发时停止。
& 粗查回答“边界在哪一段”，细查回答“边界精确到多少”；随机算法另报种子和重复次数。\\

结果及分析（p.22、p.27--28、p.32--33、p.38、p.41）
& 报告碰撞时刻、螺距、圆弧半径、曲线长度和速度上限，并解释局部异常。
& 参数微变后若第一次碰撞对象改变，论文从局部几何解释输出跳变；12--13秒的速度突增则用相同时间内的路程差解释。
& 表后不复述数字，改为回答“为何变、哪条约束起作用、结论只到什么范围”。\\

检验与边界
& 检查首次碰撞对象、全过程可行性和现实动作限制。
& p.17用反设和逐节前推说明最早碰撞对象；p.28用“终点可行但中途碰撞”的方案否定只查终点；p.32从龙身不能倒退推出半径下界。
& 检验要对应真实疑问。回代、反例和物理边界各自有用途，不能合写成一句“模型得到验证”。\\

附录（p.43--62）
& 给出逐把手循环、二分、分段路径、碰撞和粒子群代码。
& 正文四种运动状态与代码状态码1--4相互对应，表格输出也由固定变量名生成。
& 附录保留能复现正文结果的入口、参数和输出说明；正文只留关键流程，不粘贴整段程序。\\
\end{longtable}
\normalsize

\subsection{五个分问的思路怎样连起来}

\begin{enumerate}
  \item 问题一先求龙头在阿基米德螺线上的位置，再由相邻把手距离不变逐节递推；
  速度由短时间内的弧长变化得到。该问交付以后各问共用的位置和速度计算器。
  \item 问题二把每块板凳写成矩形，用分离轴判断重叠。论文没有直接缩小检查范围，
  而是先用反证法说明第一次碰撞对象，再让程序在第一次触发处停止。
  \item 问题三把“能进入调头区域”改写成全过程不发生碰撞。螺距从粗区间缩到细步长，
  终点可行却中途碰撞的候选被舍去，所以可行性不是只看最后一帧。
  \item 问题四先由相切关系写出两段圆弧，再用“后续构件不能倒退”给大圆弧半径下界；
  确定路径后，位置递推按盘入、两段圆弧和盘出四种状态继续运行。
  \item 问题五不重建运动模型，只把问题四得到的各把手速度写成上限约束，粒子群只负责
  搜索龙头速度。这样，算法承担的任务很窄，模型链的职责也随之分明。
\end{enumerate}

\subsection{最值得模仿的段落组织}

\paragraph{从局部关系写到全局运动。}
一条板凳龙看似有许多构件，但相邻把手距离固定。求得一个把手后，下一把手只需解一组
定长方程；候选根再由构件次序、运动方向和局部连续性筛选。这里的好处不是“递推法简单”，
而是正文每一步都能对应代码的一次循环。

\paragraph{先证明检查范围不会漏，再缩小程序。}
论文讨论第一次碰撞时，先假设第 $i$ 个构件更早碰撞，再由轨迹的先后关系逐节前推，
最后与龙头运动矛盾。证明承担的是“不漏掉更早事件”这一职责；完成这一步后，程序才有理由
只检查与龙头有关的候选对象。我们的论文若要剪枝，也应先给结构理由，再写计算节省。

\paragraph{用反例裁决简化方案。}
螺距约为 $42\,\mathrm{cm}$ 的方案在终点附近看似可行，但在完整运动过程中会提前碰撞。
这个反例定位了“只查终点”的缺口：简化方案在运动途中提前碰撞，
读者可据此判断模型为何需要保留时间推进。

\paragraph{异常结果回到局部轨迹。}
当速度在12--13秒出现突增时，论文比较相同时间内相邻节点走过的路程，并把差异落到小圆弧
附近的路径形状。解释顺序是异常时段、局部位置、同口径比较、速度结论。以后遇到曲线尖峰，
本队检查状态切换和约束变化，不急于将其归为数值噪声。

\subsection{本队形成的写作认识}

A053可借鉴的不是某个固定句子，而是判断出现的次序：整体对象被拆成可以计算的局部关系，
再说明数学解怎样受物理次序筛选；先证明事件缩域不会漏检，再停止后续分支；参数结果不连续时，
回到第一次触发对象和局部轨迹解释。这样的段落会留下作者实际做过的选择，也使公式、图和代码
各自承担明确任务。本队模板据此要求每一节回答“写什么、怎么写、交稿前查什么”。

\section{本队论文写作模板}

以下模板用于正式竞赛论文。标题名称保持固定；实际写作时删去本节的说明文字，正文只留下本题
需要的公式、图表、结果和论证。分问内部采用“任务--变量--模型--算法--结果--解释--检验”的顺序，
没有独立结果的环节无需单列小节。

\subsection{摘要与关键词}
\begin{writeguide}
\guideitem{写什么}{交代题目对象、每个分问的主要方法、关键数值和对应结论；末尾列3--5个能检索模型或任务的关键词。}
\guideitem{怎么写}{按分问组织短段。每段先说该问要判断什么，再写模型动作，随后给最有区分度的结果和检验口径。算法名只在它确实产生某个输出时出现。}
\guideitem{交稿前查}{摘要中的数字能否在正文表格找到；单位是否一致；是否误写了正文没有做的检验；是否只列方法而没有答案。}
\end{writeguide}

\subsection{问题重述}
\begin{writeguide}
\guideitem{写什么}{用本队语言重述背景、已知条件、各问输入和交付物，不复制大段题面。}
\guideitem{怎么写}{先说明研究对象和现实任务，再逐问写“给定什么、要求求什么”。附件字段、时间范围和提交文件在相应分问中点明。}
\guideitem{交稿前查}{有没有改变题意、遗漏硬条件或混入解题方法；变量单位和对象范围是否与题面相同。}
\end{writeguide}

\subsection{问题分析}
\begin{writeguide}
\guideitem{写什么}{说明难点、分问关系、可利用结构、拟用模型及前后问接口。}
\guideitem{怎么写}{从数据形态、物理关系或决策约束中指出真实难点，随后将任务拆成可计算步骤。若沿用前问，列出保留的输入输出和新增的变量、约束。}
\guideitem{交稿前查}{是否从“采用某算法”起笔；模型选择有没有来自题目结构；各分问是否被写成互不相干的算法清单。}
\end{writeguide}

\subsection{模型假设}
\begin{writeguide}
\guideitem{写什么}{列出题面未明确、但计算必须限定的理想化条件，并说明忽略因素。}
\guideitem{怎么写}{每条假设包含对象、条件和计算后果。例如忽略摩擦后删去哪个力项，刚体假设后哪段距离保持不变。}
\guideitem{交稿前查}{假设之间是否冲突；是否把题面事实重复成假设；某条假设若删除，模型是否完全不受影响。}
\end{writeguide}

\subsection{符号说明}
\begin{writeguide}
\guideitem{写什么}{列出跨公式反复使用的变量、参数、集合、下标和单位。}
\guideitem{怎么写}{按“符号--含义--单位”制表；向量、矩阵、随机量和决策量在排版上保持一致。仅在一处使用的中间量随公式解释，不挤进总表。}
\guideitem{交稿前查}{是否一符多义、同义多符、下标未解释或单位缺失；符号表中的量是否在正文真正出现。}
\end{writeguide}

\subsection{分问题模型建立与求解}
\begin{writeguide}
\guideitem{写什么}{每问给出输入数据、变量、方程或目标、约束、求解步骤、输出和必要图表。}
\guideitem{怎么写}{给出最简单可核的模型，并点明它在哪个可观察情形不足；只对该缺口增加项或更换方法。公式后解释量的含义，算法写清初始化、更新、可行性处理、停止条件和输出文件。}
\guideitem{交稿前查}{是否只有算法名没有数学模型；约束能否追到题意；多根、分段、随机种子、搜索范围和终止精度是否交代；前问接口是否真的可复用。}
\end{writeguide}

\subsection{结果分析与模型检验}
\begin{writeguide}
\guideitem{写什么}{报告关键结果，解释方向、数量级、约束状态和现实含义；再用与结论匹配的办法检查。}
\guideitem{怎么写}{表格给精确数值，图形展示关系或状态。正文从一个具体结果起笔，随后解释它由哪个机制或约束造成。回代、守恒、残差、交叉验证、基线对照和反例分别按实际职责命名。}
\guideitem{交稿前查}{是否只复述表中数字；图表是否在正文出现并被解释；“最优、准确、稳定”等判断有没有候选范围、指标和比较对象。}
\end{writeguide}

\subsection{灵敏度或稳健性分析}
\begin{writeguide}
\guideitem{写什么}{选择会受测量、估计或情景变化影响的关键参数，观察输出数值与决策是否改变。}
\guideitem{怎么写}{先给基准值、扰动来源和范围；每次写清固定量、改变量和输出。连续结果报变化幅度，离散决策说明何时换方案。}
\guideitem{交稿前查}{是否把求解步长加密误当成灵敏度；不同参数是否用了可比较的扰动尺度；是否只展示最好看的方向。}
\end{writeguide}

\subsection{模型评价与改进}
\begin{writeguide}
\guideitem{写什么}{评价本模型能够回答的任务、计算和解释方面的长处，以及尚未纳入的因素和对应修改办法。}
\guideitem{怎么写}{优点回到公式、数据接口或检验结果；不足指出受影响的对象和边界。改进项写成“新增什么数据或关系、改哪一层、预期检查什么”，不写空泛展望。}
\guideitem{交稿前查}{是否把速度快、精度高、适用面广当成无证据套话；改进是否超出当前数据，是否被误写成现有成果。}
\end{writeguide}

\subsection{参考文献}
\begin{writeguide}
\guideitem{写什么}{列出正文实际引用的题目资料、方法来源、数据说明和软件文档。}
\guideitem{怎么写}{正文在首次使用外部定义、算法或数据时给引用；文末按统一格式列作者、题名、载体、年份等现有信息。}
\guideitem{交稿前查}{是否有文后未引用条目或正文无出处的材料；编号、年份、链接和访问日期是否完整；是否把搜索结果页当作来源。}
\end{writeguide}

\subsection{附录}
\begin{writeguide}
\guideitem{写什么}{给出主要代码入口、参数表、完整中间结果、补充推导和题目要求的输出文件说明。}
\guideitem{怎么写}{先列文件清单和运行顺序，再给关键函数；代码变量与正文符号建立对照，随机算法写种子和预算。正文放裁决所需结果，超长表放附录。}
\guideitem{交稿前查}{从原始附件到关键表能否按说明运行；是否缺依赖、参数、单位或输出路径；正文与附录的变量名和数值是否一致。}
\end{writeguide}

\appendix
\section{模板提交前速查}

\begin{enumerate}
  \item 摘要中的结论、正文表格和附件输出使用同一数值与单位；
  \item 每个模型都有变量、方程或目标、约束、求解入口和结果解释；
  \item 每幅图和每张表在正文中被提出，并至少承担一次读数或比较；
  \item 关键结论配有回代、误差、交叉验证、基线、反例或物理边界中的一种检查；
  \item 附录列出运行顺序、依赖、随机种子、输入文件和输出文件。
\end{enumerate}

\section*{材料来源说明}
\addcontentsline{toc}{section}{材料来源说明}

\begin{thebibliography}{9}
\bibitem{mindmap}
数学建模学习思维导图，工作区文件 \texttt{1.png}，2026年7月核读。

\bibitem{A053}
编号A053：\emph{基于目标优化的板凳龙形态调节}，2024年全国大学生数学建模竞赛A题一等奖论文语料，
工作区文件 \texttt{2024/A053.pdf}。本文所列数值按论文页面记录，附录代码已阅读，未在本报告中独立运行MATLAB复算。

\bibitem{problem2024A}
2024年全国大学生数学建模竞赛A题：板凳龙，工作区赛题原文。
\end{thebibliography}

\end{document}
